\chapter{Skąd bierze się 249\,219 parametrów}\label{dod:parametry}

Tabela rozpisuje każdy moduł prawdziwego, wytrenowanego modelu Gen~6 ($H=48$, $D=64$, $n_h=4$, $S=53$, $W=43$). Każdą liczbę da się policzyć ze wzorów: warstwa liniowa $n\to m$ ma $mn+m$ parametrów (rozdział~\ref{rozdz:wektory}),
konwolucja $C_\mathrm{we}C_\mathrm{wy}k_tk_p+C_\mathrm{wy}$ (tw.~\ref{tw:param-konw}), LayerNorm $2\times$ wymiar. Sumy zostały porównane programem z liczbą parametrów modelu wczytanego z pliku \texttt{best.pt}.

\begin{table}[ht]
\centering\small
\begin{tabular}{@{}p{5.4cm}p{3.7cm}r@{}}
\toprule
moduł & obliczenie & parametry\\
\midrule
trzon: konwolucja $3{\times}7$, $9\to48$ & $9\cdot48\cdot21+48$ & $9\,120$\\
\emph{blok harmoniczny} ($\times3$): LN$(48)$ & $2\cdot48$ & $96$\\
\quad konwolucja $1{\times}1$, $336\to48$ & $336\cdot48+48$ & $16\,176$\\
\quad konwolucja głębokościowa $1{\times}5$ & $48\cdot5+48$ & $288$\\
\quad konwolucja czasowa $3{\times}1$, $48\to48$ & $48\cdot48\cdot3+48$ & $6\,960$\\
\quad suma bloku $\times3$ & $23\,520\cdot3$ & $70\,560$\\
zejście: konwolucja $1{\times}3$, $48\to64$ & $48\cdot64\cdot3+64$ & $9\,280$\\
osadzenie wysokości $e_s$ & $53\cdot64$ & $3\,392$\\
\emph{warstwa kontekstu} ($\times2$): trzy LN & $3\cdot128$ & $384$\\
\quad uwaga w czasie ($W_{QKV},W_O$, premie) & $12\,480+4160+4\cdot87$ & $16\,988$\\
\quad uwaga między wysokościami & $12\,480+4160+4\cdot105$ & $17\,060$\\
\quad FF: $64\to256\to64$ & $16\,640+16\,448$ & $33\,088$\\
\quad suma warstwy $\times2$ & $67\,520\cdot2$ & $135\,040$\\
LN końcowy & $2\cdot64$ & $128$\\
głowica początku & $(64\cdot64+64)+(64+1)$ & $4\,225$\\
głowica „brzmi” ($65$ wejść) & $(65\cdot64+64)+65$ & $4\,289$\\
głowica końca & jak początku & $4\,225$\\
głowica głośności (nieuczona) & jak początku & $4\,225$\\
głowica tabulatury: $\vect w_i,c_i$; premie progów; MLP ciszy & $390+120+4\,225$ & $4\,735$\\
\midrule
\textbf{razem} & & $\mathbf{249\,219}$\\
\bottomrule
\end{tabular}
\caption{Rozbicie liczby parametrów Gen~6 (zgodne z tw.~\ref{tw:liczba-parametrow}).}
\label{tab:parametry-pelne}
\end{table}

\begin{przyklad}[Kontrola sumy]\label{prz:suma-kontrola}
$9120+70\,560+9280+3392=92\,352$ (pierwsza połowa), $+\,135\,040$ (kontekst) $=227\,392$, $+\,128+4225+4289+4225+4225+4735=21\,827$ (głowice i LN), razem $249\,219$.
\end{przyklad}
